\subsection{Overall Architecture}
\label{sec:overall-architecture}

Breakly is organised as a mobile-first system. The mobile client contains the
presentation layer, break-session orchestration, local persistence, and
notification integration. An optional backend can provide authenticated
content synchronisation and analytics without being required for the core
offline flow. External services are isolated behind adapters so that the
application remains testable when a service is unavailable.

\begin{figure}[H]
\centering
\fbox{\parbox{0.85\textwidth}{\centering
Mobile app $\leftrightarrow$ API/backend $\leftrightarrow$ database and
content services\\[4pt]
The mobile app also uses local storage and the operating-system notification
service directly.}}
\caption{Breakly overall architecture placeholder}
\label{fig:overall-architecture}
\end{figure}

\subsection{In-App Architecture}
\label{sec:in-app-architecture}

The app follows a pragmatic Clean Architecture with MVVM at the presentation
boundary:
\begin{itemize}
\item \textbf{Presentation:} screens, navigation, view models, and reusable
  UI components.
\item \textbf{Domain:} break-session rules, duration validation, activity
  selection, and session result policies.
\item \textbf{Data:} repositories, local storage, optional API client, and
  notification adapter.
\end{itemize}

This separation keeps timer and preference rules independent of React
components and makes them unit-testable. React Native with Expo is selected
for cross-platform delivery, TypeScript for type safety, and a local
database or key-value store for offline-first persistence. Final package
versions and the decision matrix will be recorded before implementation.

\subsection{Main Flow Sequences}
\label{sec:sequence-diagrams}

The final report will include sequence diagrams for: starting a break,
pausing/resuming and completing a break, scheduling a reminder, and loading
content while offline. Each diagram must show the user, screen/view model,
domain service, repository, storage, and notification service where relevant,
including failure paths.
